%=============================================================================!
% 5.S  WHAT WE NOW HAVE                                                       !
%=============================================================================!
\unnumbered{What we now have}

A body turning about a point has angular momentum $\vb{L} = \vb{r}\times
\vb{p}$, which changes at the rate of the torque $\vb{G} = \vb{r}\times
\vb{F}$; in a plane these are $L = mrv_\theta$ and $G = rF_\theta$, the
lever. A central force exerts no torque, so it conserves $\vb{L}$: a puck
pulled in on a string goes round faster, sweeping equal areas in equal
times. For a group of bodies whose forces act along the lines joining
them, as forces between atoms built from pairs do, the internal torques
cancel and only external torques change the total, which splits into a part
carried by the centre of mass and a part about it.

A rigid body turning
with angular velocity $\boldsymbol{\omega}$ moves with $\vb{v} =
\boldsymbol{\omega}\times\vb{r}$; about a fixed axis $L_z = I\omega$ and $K
= \frac12 I\omega^2$, with $I = m_{\mathrm r}d^2$ for two bodies, and a
rolling body keeps part of its energy in turning.

In general $\vb{L} =
\mathsf{I}\boldsymbol{\omega}$ and $K = \frac12\boldsymbol{\omega}\cdot
\mathsf{I}\boldsymbol{\omega}$ with the symmetric inertia tensor, whose
eigenvectors are the principal axes. Subtracting the drift $\vb{V} =
\vb{P}/M$ and the turning $\boldsymbol{\omega}\times\vb{r}'$ with
$\mathsf{I}\boldsymbol{\omega} = \vb{L}'$ removes the momentum and
the angular momentum from a set of velocities and splits the kinetic energy
into three separate parts.

The rigid motions are zero modes of the
Hessian, which leaves a molecule $3N - 6$ ways to vibrate ($3N - 5$ if it
is linear), and at room temperature a molecule such as LiF turns once in
the time it vibrates 27 times.

Chapter~6 rebuilds mechanics from the energy alone, in a form that
handles constraints, such as the rod of a pendulum, without ever finding
the forces they exert.
